\section{四个参数详解}

售后日常只需要调四个参数：等待时间、运行时间、加速度、速度。这一章逐个讲清楚每个参数"是什么""背后在调节什么""怎么调"。

\textbf{先理解一个前提：}这套设备是一台自动折盒机，传送带是整套流水线的\textbf{最后一道工序}。当前面折纸机构完成折纸后，传送带把折好的纸盒送出去。四个参数的核心目的不是"控制传送多远"——传送距离由机械结构固定——而是\textbf{让传送带这道工序和前后工序在时间上配合好}。这就是"时序协调"。

\subsection{等待时间}

\textbf{是什么：}前一道工序（折纸机构）完成后，延迟多久才启动传送带。

\textbf{背后在调节什么：}

当前面折纸机构完成折纸动作、下方吸盘停止吸气后，纸盒不再被吸附固定——这个"吸盘停止吸气"的信号就是传送带模块的外部启动信号。

但吸盘刚停止吸气时，纸盒可能还没有完全稳定。等待时间就是在"前工序结束"和"传送带启动"之间插入的一段缓冲，让纸盒稳定释放后再启动传送带。本质上是和前道工序的\textbf{时序配合}。

你可以把它理解为：前一道工序说"我干完了"，等待时间就是"收到通知后倒数几秒再起跑"。倒数太短，可能前工序还没彻底收好你就启动了；倒数太长，前后工序之间就会出现空档，拖慢整体节拍。

\begin{itemize}
    \item \textbf{默认值：}\texttt{250 ms}
    \item \textbf{可调范围：}\texttt{0 \textasciitilde{} 10000 ms}
    \item \textbf{怎么调：}
    \begin{itemize}
        \item 传送带启动太早，和前工序冲突$\rightarrow$ \textbf{增大}等待时间
        \item 传送带启动太晚，前后工序有空档$\rightarrow$ \textbf{减小}等待时间
        \item 希望前工序一结束马上启动$\rightarrow$ 设为 \texttt{0}
    \end{itemize}
\end{itemize}

改完后点"保存"。

\begin{figure}[H]
    \centering
    \includegraphics[width=0.7\textwidth]{figures/ui-runtime-params.png}
    \caption{运行参数区域。运行时间和等待时间的滑块与保存按钮。}
    \label{fig:ui-runtime-params}
\end{figure}

\subsection{运行时间}

\textbf{是什么：}传送带电机在一个周期内持续运行多久。

\textbf{背后在调节什么：}

在这台折盒机里，传送带的传送距离是由机械结构固定的——跑道的长度不会变。运行时间配合速度，共同决定了传送带完成一次送料需要多久，以及何时准备好接收下一个纸盒。

运行时间过短，传送带还没把纸盒送到指定位置就停了；运行时间过长，传送带空转浪费时间，可能和下一周期的前工序发生冲突。本质上是在调节本道工序的"工作时长"，让它和整套流水线的\textbf{节拍匹配}。

\begin{itemize}
    \item \textbf{默认值：}\texttt{250 ms}
    \item \textbf{可调范围：}\texttt{1 \textasciitilde{} 60000 ms}
    \item \textbf{怎么调：}
    \begin{itemize}
        \item 纸盒没送到指定位置（运行时间太短）$\rightarrow$ \textbf{增大}运行时间
        \item 运行时间太长，和下一周期冲突$\rightarrow$ \textbf{减小}运行时间
    \end{itemize}
\end{itemize}

改完后点"保存"。

\subsection{加速度}

\textbf{是什么：}电机启动时，从静止加速到目标速度的快慢程度。界面显示 \texttt{0 \textasciitilde{} 100}，数值越大，启动越快。

\textbf{背后在调节什么：}

电机从静止瞬间跳到全速，会产生很大的机械冲击，长期下去可能损坏传送机构或让物料移位。加速度参数控制的就是这个"起步过程有多柔和"。

这里需要理解一点：界面上显示的"加速度"和固件底层的"加速时间"大致对应，但不是严格的一一换算。

\begin{itemize}
    \item 界面显示 \texttt{100} $\rightarrow$ 固件加速时间约为 \texttt{0 ms}，意味着几乎没有加速过程，直接全速启动
    \item 界面显示 \texttt{0} $\rightarrow$ 固件加速时间约为 \texttt{10000 ms}，意味着用 10 秒钟慢慢爬到目标速度
\end{itemize}

所以售后只需要记住：\textbf{数值越大，启动越冲；数值越小，启动越柔和。}

\begin{itemize}
    \item \textbf{默认值：}\texttt{100}（无加速，直接全速启动）
    \item \textbf{可调范围：}\texttt{0 \textasciitilde{} 100}
    \item \textbf{怎么调：}
    \begin{itemize}
        \item 启动瞬间冲击明显（机械振动大、噪音大）$\rightarrow$ \textbf{减小}加速度，例如从 \texttt{100} 调到 \texttt{70} 或 \texttt{60}
        \item 启动太慢、拖得太久才达到正常速度 $\rightarrow$ \textbf{增大}加速度
    \end{itemize}
\end{itemize}

改完后点"保存"。

\begin{figure}[H]
    \centering
    \includegraphics[width=0.7\textwidth]{figures/ui-speed-params.png}
    \caption{速度参数区域。速度（应用/保存双按钮）、加速度（保存按钮），均带滑块。}
    \label{fig:ui-speed-params}
\end{figure}

\subsection{速度}

\textbf{是什么：}传送带正常运行时的速度，单位是厘米/分钟（cm/min）。数值越大，跑得越快。

\textbf{背后在调节什么：}

这台折盒机的传送距离是固定的（由机械结构决定），所以速度本质上控制的是\textbf{"完成一次送料需要多久"}——也就是本道工序的节拍。

\begin{itemize}
    \item 速度快 $\rightarrow$ 完成送料快 $\rightarrow$ 本道工序节拍快
    \item 速度慢 $\rightarrow$ 完成送料慢 $\rightarrow$ 本道工序节拍慢
\end{itemize}

如果本道工序节拍和前面折纸机构的节拍不匹配，就会出现"前工序已经送下一个纸盒来了，但传送带还没把上一个送走"的堵塞，或者"传送带空等"的效率损失。

界面上显示的速度值，是上位机帮你换算好的直观数值。底层用的是 PWM（脉宽调制）来控制电机转速，但售后不需要直接操作 PWM——调"速度"更直观。

默认速度 \texttt{48226 cm/min} 是从底层默认 PWM \texttt{150} 和速度校准值 \texttt{3655 RPM} 换算得到的。

\begin{itemize}
    \item \textbf{默认值：}\texttt{48226 cm/min}
    \item \textbf{可调范围：}\texttt{0 \textasciitilde{} 70000 cm/min}
\end{itemize}

\begin{quote}
\textbf{注意：}速度设为 \texttt{0} 时，电机不会转动。如果点了"单次运行测试"但设备没动作，先检查一下速度是不是被误设成了 \texttt{0}。
\end{quote}

\textbf{操作方式（重点）：}

速度参数有两个按钮，用法和其他参数不一样：

\begin{table}[H]
    \centering
    \caption{速度参数按钮说明}
    \label{tab:speed-buttons}
    \begin{tabular}{lcp{6cm}}
        \toprule
        \textbf{按钮} & \textbf{作用} & \textbf{断电后是否保留} \\
        \midrule
        应用 & 临时试一下当前速度 & 不保留 \\
        保存 & 保存为设备默认速度 & 保留 \\
        \bottomrule
    \end{tabular}
\end{table}

\textbf{推荐用法：}

\begin{enumerate}[label={(\arabic*)}]
    \item 在速度输入框里填入想试的数值。
    \item 点"应用"。此时速度临时生效，但还没有写进设备。
    \item 点"单次运行测试"，观察效果。
    \item 速度合适后，点"保存"。保存成功后才真正写入设备，断电后仍然保留。
\end{enumerate}

\textbf{怎么调：}
\begin{itemize}
    \item 整体节拍太慢，前工序等不及$\rightarrow$ \textbf{增大}速度
    \item 送料过快/不稳$\rightarrow$ \textbf{减小}速度
\end{itemize}
